\chapter{Writing}
\label{ch:writing}

This book has drawn most of its examples from source code and structured data, consistent with Vim's reputation as a programmer's editor, but nothing in the grammar covered throughout Parts II through VI is specific to code. This closing chapter of Part VII turns to long-form prose: Markdown, LaTeX, and the particular demands of a manuscript large enough to span many files and many editing sessions, the kind of writing this book itself was produced with.

\section{Markdown}

Markdown's structural elements, headings, emphasis markers, links, and code spans, are plain text with light punctuation conventions, which means they are addressable with the ordinary motions and text objects of Part II rather than requiring any Markdown-specific commands. A heading is a line beginning with one or more \texttt{\#} characters, navigable with the paragraph and section motions of Chapter~\ref{ch:coordinates} depending on how a document is structured; emphasis markers (\texttt{*}, \texttt{**}, \texttt{\_}) delimit a span exactly as quotes do, and while Vim has no dedicated built-in text object for "the word between asterisks," the quote text objects from Chapter~\ref{ch:regions}, \texttt{i*} is not built in, though a reader who finds themselves wanting it is precisely the audience for the user-defined text objects noted at the end of that chapter.

The formatting operators \texttt{gq} and \texttt{gw} from Chapter~\ref{ch:operators} are particularly suited to Markdown prose, rewrapping a paragraph to the configured \texttt{textwidth} after an edit has left its line lengths ragged; setting \texttt{textwidth=80} (or whatever width a particular publishing target expects) and reflowing edited paragraphs with \texttt{gqip} as a matter of habit keeps a Markdown document's source consistently formatted even under heavy revision, which matters more than it might first appear, since a version-controlled document with inconsistent line wrapping produces much noisier diffs than one with consistent wrapping, obscuring genuine content changes among incidental rewrapping in every diff review.

\section{LaTeX}

LaTeX introduces genuine structural complexity beyond Markdown's light punctuation, and several of Vim's built-in facilities apply to it directly and usefully. The tag text object from Chapter~\ref{ch:regions}, \texttt{it}/\texttt{at}, is defined for markup with paired opening and closing delimiters and does not, out of the box, recognize LaTeX's \texttt{\textbackslash begin\{...\}} / \texttt{\textbackslash end\{...\}} environment syntax the way it recognizes HTML or XML tags; a reader working in LaTeX heavily enough to want an equivalent should expect to reach for a filetype-specific plugin or a custom text object rather than the built-in \texttt{it}/\texttt{at} directly, consistent with this book's plugin-free scope simply noting where a natural extension point exists rather than claiming built-in coverage that isn't there. Percent-delimited comments, section and subsection commands, and citation keys, by contrast, are all addressable with the ordinary word and WORD motions of Chapter~\ref{ch:coordinates} without any special handling at all, since they are simply punctuation-delimited text like any other.

\texttt{:set spell} enables spell-checking directly in the buffer, underlining suspect words without disrupting the surrounding LaTeX markup, and \texttt{]s} and \texttt{[s} jump to the next and previous flagged word; \texttt{z=} at a flagged word suggests corrections, and \texttt{zg} adds a word to the session's known-good word list, useful for the field-specific terminology and citation keys that inevitably appear throughout a technical manuscript and would otherwise be flagged repeatedly.

\section{Large Manuscripts}

A book-length manuscript, split across many files (one per chapter, in the manner this very book's own source is organized), benefits from several of the multi-file mechanisms covered earlier in this book, brought together here for the specific case of long-form writing. The argument list from Chapter~\ref{ch:global}, populated with \texttt{:args chapters/*.tex} or a similar glob, makes \texttt{:argdo} available across every chapter file at once: a terminology change, a formatting correction, or any of the substitution patterns from Chapter~\ref{ch:substitution} can be applied uniformly across an entire manuscript in a single command, exactly as Chapter~\ref{ch:global} demonstrated for smaller-scale editing tasks.

\texttt{:vimgrep} \textit{pattern} \texttt{chapters/*.tex}, from Chapter~\ref{ch:workflows}, searches across every chapter file for a pattern, populating the quickfix list with every match, which is the practical way to answer a question like "which chapters mention this term" or "where did I leave that unfinished cross-reference" across a manuscript too large to hold in memory or search file by file. Global marks from Chapter~\ref{ch:marks} are worth setting deliberately at points a writer expects to return to repeatedly, such as an outline file or a running notes document, since they remain addressable from within any chapter file without first navigating there manually.

Finally, a session, from Chapter~\ref{ch:buffers}, is the natural unit for a long-running writing project specifically: saving the exact arrangement of open chapter files, splits, and cursor positions at the end of a working session and restoring it at the start of the next removes the friction of re-opening and re-navigating to the same working set of files every time work resumes, which matters more for a manuscript revisited over weeks or months than for almost any other kind of editing task covered in this book.

\section{Notes and Documentation}

Plain-text notes, whether a single running file or a directory of many small ones, benefit from a lighter version of the same large-manuscript machinery. Marks and the jump list from Chapter~\ref{ch:marks} suffice for a single large notes file; for a directory of many small ones, \texttt{:e} combined with the filename completion from Chapter~\ref{ch:ex} is frequently faster than any more elaborate navigation scheme, particularly once note filenames follow a consistent, predictable convention that completion can exploit. The global command from Chapter~\ref{ch:global}, applied across a set of notes opened via the argument list, answers questions like "which notes mention this project" exactly as it would across a manuscript's chapters, since a directory of notes and a book's chapter files are, from Vim's perspective, the same kind of thing: a set of plain-text buffers addressable together.

\section{Writing as the Terminus of This Book's Argument}

This chapter closes Part VII, and with it every applied chapter this book has to offer before Part VIII's reference material begins. It is worth noting, in closing, what has and has not been true throughout this chapter: not one command introduced here is new. Every mechanism this chapter has applied to Markdown, LaTeX, and large manuscripts was already established, for other purposes, somewhere in Parts II through VI. This is not a coincidence and not a limitation; it is the entire thesis of this book, restated one final time in its most human-facing context. A reader who arrived at this chapter already fluent in the grammar covered earlier did not need a separate vocabulary for writing prose, because Vim does not have one. Text is text, whether it is a function definition, a rebase plan, a merge conflict, or a paragraph of English prose, and the same small set of primitives, composed the same small number of ways, was sufficient for all of it.
